% Gemeinsame Originalaussagen; nur Band 45 registriert sie.
\newcommand{\SemilatticeStatement}[1]{\csname SemilatticeStatement@#1\endcsname}

\expandafter\newcommand\csname SemilatticeStatement@JoinOrderFirstCarrier\endcsname{%
\FormulaThmDeltaK[Erste Komponente der supremal induzierten Relation]{%
  \Semi{A}\dsep x\InducedLeq y\vdash x\in A
}{JoinOrderFirstCarrier}{%
  \DeltaRow{Mengen}{A\dsep x\dsep y}
  \DeltaPrem{Von \(\Semi{A}\) supremal induzierter Relationsoperator}{\InducedLeq}
}
}

\expandafter\newcommand\csname SemilatticeStatement@JoinOrderSecondCarrier\endcsname{%
\FormulaThmDeltaK[Zweite Komponente der supremal induzierten Relation]{%
  \Semi{A}\dsep x\InducedLeq y\vdash y\in A
}{JoinOrderSecondCarrier}{%
  \DeltaRow{Mengen}{A\dsep x\dsep y}
  \DeltaPrem{Von \(\Semi{A}\) supremal induzierter Relationsoperator}{\InducedLeq}
}
}

\expandafter\newcommand\csname SemilatticeStatement@JoinOrderEvaluation\endcsname{%
\FormulaThmDeltaK[Auswertung der supremal induzierten Relation]{%
  \Semi{A}\dsep x\InducedLeq y
  \vdash x\OmittedSemiOp{}y=y
}{JoinOrderEvaluation}{%
  \DeltaRow{Mengen}{A\dsep x\dsep y}
  \DeltaPrem{Von \(\Semi{A}\) supremal induzierter Relationsoperator}{\InducedLeq}
}
}

\expandafter\newcommand\csname SemilatticeStatement@JoinOrderReflexive\endcsname{%
\FormulaThmDeltaK[Reflexivität der supremal induzierten Relation]{%
  \Semi{A}\dsep x\in A\vdash x\InducedLeq x
}{JoinOrderReflexive}{%
  \DeltaRow{Mengen}{A\dsep x}
  \DeltaPrem{Von \(\Semi{A}\) supremal induzierter Relationsoperator}{\InducedLeq}
}
}

\expandafter\newcommand\csname SemilatticeStatement@JoinOrderTransitive\endcsname{%
\FormulaThmDeltaK[Transitivität der supremal induzierten Relation]{%
  \Semi{A}\dsep x\InducedLeq y
  \dsep y\InducedLeq z\vdash x\InducedLeq z
}{JoinOrderTransitive}{%
  \DeltaRow{Mengen}{A\dsep x\dsep y\dsep z}
  \DeltaPrem{Von \(\Semi{A}\) supremal induzierter Relationsoperator}{\InducedLeq}
}
}

\expandafter\newcommand\csname SemilatticeStatement@JoinOrderAntisymmetric\endcsname{%
\FormulaThmDeltaK[Antisymmetrie der supremal induzierten Relation]{%
  \Semi{A}\dsep x\InducedLeq y
  \dsep y\InducedLeq x\vdash x=y
}{JoinOrderAntisymmetric}{%
  \DeltaRow{Mengen}{A\dsep x\dsep y}
  \DeltaPrem{Von \(\Semi{A}\) supremal induzierter Relationsoperator}{\InducedLeq}
}
}

\expandafter\newcommand\csname SemilatticeStatement@JoinSemilatticeInducesOrder\endcsname{%
\FormulaThmDeltaK[Ein supremal gelesener Halbverband induziert eine partielle Ordnung]{%
  \Semi{A}\vdash \PartOrd{A,\InducedLeq}
}{JoinSemilatticeInducesOrder}{%
  \DeltaRow{Mengen}{A\dsep x\dsep y\dsep z}
  \DeltaPrem{Von \(\Semi{A}\) supremal induzierter Relationsoperator}{\InducedLeq}
}
}

\expandafter\newcommand\csname SemilatticeStatement@JoinUpperLeft\endcsname{%
\FormulaThmDeltaK[Der erste Operand liegt unter dem algebraischen Supremum]{%
  \Semi{A}\dsep x,y\in A
  \vdash x\InducedLeq x\OmittedSemiOp{}y
}{JoinUpperLeft}{%
  \DeltaRow{Mengen}{A\dsep x\dsep y}
  \DeltaPrem{Von \(\Semi{A}\) supremal induzierter Relationsoperator}{\InducedLeq}
}
}

\expandafter\newcommand\csname SemilatticeStatement@JoinUpperRight\endcsname{%
\FormulaThmDeltaK[Der zweite Operand liegt unter dem algebraischen Supremum]{%
  \Semi{A}\dsep x,y\in A
  \vdash y\InducedLeq x\OmittedSemiOp{}y
}{JoinUpperRight}{%
  \DeltaRow{Mengen}{A\dsep x\dsep y}
  \DeltaPrem{Von \(\Semi{A}\) supremal induzierter Relationsoperator}{\InducedLeq}
}
}

\expandafter\newcommand\csname SemilatticeStatement@JoinLeastUpper\endcsname{%
\FormulaThmDeltaK[Das algebraische Supremum ist die kleinste obere Schranke]{%
  \begin{aligned}[t]
    &\Semi{A}\dsep x,y,u\in A\dsep
      x\InducedLeq u\dsep y\InducedLeq u\vdash{}\\
    &x\OmittedSemiOp{}y\InducedLeq u
  \end{aligned}
}{JoinLeastUpper}{%
  \DeltaRow{Mengen}{A\dsep x\dsep y\dsep u}
  \DeltaPrem{Von \(\Semi{A}\) supremal induzierter Relationsoperator}{\InducedLeq}
}
}

\expandafter\newcommand\csname SemilatticeStatement@JoinIsPairSupremum\endcsname{%
\FormulaThmDeltaK[Die Halbverbandsoperation liefert das Paarsupremum]{%
  \Semi{A}\dsep x,y\in A
  \vdash x\OmittedSemiOp{}y=\sup_{A,\InducedLeq}(\{x,y\})
}{JoinIsPairSupremum}{%
  \DeltaRow{Mengen}{A\dsep x\dsep y\dsep u\dsep z}
  \DeltaPrem{Von \(\Semi{A}\) supremal induzierter Relationsoperator}{\InducedLeq}
}
}

\expandafter\newcommand\csname SemilatticeStatement@PairJoinUpperLeft\endcsname{%
\FormulaThmDeltaK[Erster Operand einer Paarsupremumsoperation]{%
  \begin{aligned}[t]
    &\PartOrd{A,\leq}\dsep\PairJoinOp{A,\NeutralSemiOp}
      \dsep x,y\in A\vdash{}\\
    &x\leq x\NeutralSemiOp{}y
  \end{aligned}
}{PairJoinUpperLeft}{%
  \DeltaRow{Mengen}{A\dsep x\dsep y\dsep m\dsep u}
  \DeltaRow{Relationsoperatoren}{\leq}
  \DeltaRow{Funktionen}{\NeutralSemiOp}
}
}

\expandafter\newcommand\csname SemilatticeStatement@PairJoinUpperRight\endcsname{%
\FormulaThmDeltaK[Zweiter Operand einer Paarsupremumsoperation]{%
  \PartOrd{A,\leq}\dsep\PairJoinOp{A,\NeutralSemiOp}\dsep x,y\in A
  \vdash y\leq x\NeutralSemiOp{}y
}{PairJoinUpperRight}{%
  \DeltaRow{Mengen}{A\dsep x\dsep y}
  \DeltaRow{Relationsoperatoren}{\leq}
  \DeltaRow{Funktionen}{\NeutralSemiOp}
}
}

\expandafter\newcommand\csname SemilatticeStatement@PairJoinLeastUpper\endcsname{%
\FormulaThmDeltaK[Kleinste obere Schranke einer Paarsupremumsoperation]{%
  \begin{aligned}[t]
    &\PartOrd{A,\leq}\dsep\PairJoinOp{A,\NeutralSemiOp}
      \dsep x,y,u\in A\dsep x\leq u\dsep y\leq u\vdash{}\\
    &x\NeutralSemiOp{}y\leq u
  \end{aligned}
}{PairJoinLeastUpper}{%
  \DeltaRow{Mengen}{A\dsep x\dsep y\dsep u\dsep z\dsep m}
  \DeltaRow{Relationsoperatoren}{\leq}
  \DeltaRow{Funktionen}{\NeutralSemiOp}
}
}

\expandafter\newcommand\csname SemilatticeStatement@PairJoinIdempotent\endcsname{%
\FormulaThmDeltaK[Idempotenz einer Paarsupremumsoperation]{%
  \PartOrd{A,\leq}\dsep \PairJoinOp{A,\NeutralSemiOp}\dsep x\in A
  \vdash x\NeutralSemiOp{}x=x
}{PairJoinIdempotent}{%
  \DeltaRow{Mengen}{A\dsep x}
  \DeltaRow{Relationsoperatoren}{\leq}
  \DeltaRow{Funktionen}{\NeutralSemiOp}
}
}

\expandafter\newcommand\csname SemilatticeStatement@PairJoinCommutative\endcsname{%
\FormulaThmDeltaK[Kommutativität einer Paarsupremumsoperation]{%
  \PartOrd{A,\leq}\dsep \PairJoinOp{A,\NeutralSemiOp}\dsep x,y\in A
  \vdash x\NeutralSemiOp{}y=y\NeutralSemiOp{}x
}{PairJoinCommutative}{%
  \DeltaRow{Mengen}{A\dsep x\dsep y}
  \DeltaRow{Relationsoperatoren}{\leq}
  \DeltaRow{Funktionen}{\NeutralSemiOp}
}
}

\expandafter\newcommand\csname SemilatticeStatement@PairJoinAssociative\endcsname{%
\FormulaThmDeltaK[Assoziativität einer Paarsupremumsoperation]{%
  \begin{aligned}[t]
    &\PartOrd{A,\leq}\dsep \PairJoinOp{A,\NeutralSemiOp}
      \dsep x,y,z\in A\vdash{}\\
    &\bigl(x\NeutralSemiOp{}y\bigr)\NeutralSemiOp{}z
      =x\NeutralSemiOp\bigl(y\NeutralSemiOp{}z\bigr)
  \end{aligned}
}{PairJoinAssociative}{%
  \DeltaRow{Mengen}{A\dsep x\dsep y\dsep z}
  \DeltaRow{Relationsoperatoren}{\leq}
  \DeltaRow{Funktionen}{\NeutralSemiOp}
}
}

\expandafter\newcommand\csname SemilatticeStatement@PairJoinCharacterization\endcsname{%
\FormulaThmDeltaK[Charakterisierung durch Paarsuprema]{%
  \PartOrd{A,\leq}\dsep \PairJoinOp{A,\NeutralSemiOp}
  \vdash \Semi{A,\NeutralSemiOp}
}{PairJoinCharacterization}{%
  \DeltaRow{Mengen}{A\dsep x\dsep y\dsep z}
  \DeltaRow{Relationsoperatoren}{\leq}
  \DeltaRow{Funktionen}{\NeutralSemiOp}
}
}

\expandafter\newcommand\csname SemilatticeStatement@PairJoinRecoversOrder\endcsname{%
\FormulaThmDeltaK[Die Paarsupremumsoperation gewinnt die Ordnung zurück]{%
  \begin{aligned}[t]
    &\PartOrd{A,\leq}\dsep \PairJoinOp{A,\NeutralSemiOp}
      \dsep x,y\in A\vdash{}\\
    &x\leq y\leftrightarrow x\InducedLeq y
  \end{aligned}
}{PairJoinRecoversOrder}{%
  \DeltaRow{Mengen}{A\dsep x\dsep y}
  \DeltaRow{Relationsoperatoren}{\leq}
  \DeltaRow{Funktionen}{\NeutralSemiOp}
  \DeltaPrem{Von \(\Semi{A,\NeutralSemiOp}\) supremal induziert}{\InducedLeq}
}
}

\expandafter\newcommand\csname SemilatticeStatement@PowerSetUnionClosure\endcsname{%
\FormulaThmDeltaK[Abgeschlossenheit der Vereinigung auf einer Potenzmenge]{%
  X,Y\in\powerset(S)\vdash X\cup Y\in\powerset(S)
}{PowerSetUnionClosure}{%
  \DeltaRow{Mengen}{S\dsep X\dsep Y}
}
}

\expandafter\newcommand\csname SemilatticeStatement@PowerSetUnionSemilattice\endcsname{%
\FormulaThmDeltaK[Vereinigung ist eine Halbverbandsoperation]{%
  \Semi{\powerset(S),\cup}
}{PowerSetUnionSemilattice}{%
  \DeltaRow{Mengen}{S\dsep X\dsep Y\dsep Z}
  \DeltaRow{Funktionen}{\cup}
}
}

\expandafter\newcommand\csname SemilatticeStatement@SubsetIffUnionEqualsRight\endcsname{%
\FormulaThmDeltaK[Inklusion durch Vereinigung]{%
  X,Y\in\powerset(S)\vdash
  \bigl(X\subseteq Y\leftrightarrow X\cup Y=Y\bigr)
}{SubsetIffUnionEqualsRight}{%
  \DeltaRow{Mengen}{S\dsep X\dsep Y}
}
}

\expandafter\newcommand\csname SemilatticeStatement@PowerSetUnionOrderIsInclusion\endcsname{%
\FormulaThmDeltaK[Die Vereinigungsrelation ist die Inklusion]{%
  X,Y\in\powerset(S)\vdash
  \bigl(X\InducedLeq Y\leftrightarrow X\subseteq Y\bigr)
}{PowerSetUnionOrderIsInclusion}{%
  \DeltaRow{Mengen}{S\dsep X\dsep Y}
  \DeltaPrem{\makecell[l]{Supremal durch
    \(\Semi{\powerset(S),\cup}\)\\induzierter Relationsoperator}}{\InducedLeq}
}
}
